% !TEX TS-program = pdflatex
% !TeX program = pdflatex
% !TEX encoding = UTF-8
% !TeX spellcheck = en_US

\documentclass[xcolor=table]{beamer}

\usepackage{../extra/beamer/karimml}

\input{options2}

\subtitle[Attention]{Attention}  

\changegraphpath{../img/attention/}

\begin{document}
	
	\begin{frame}
		\frametitle{\inserttitle}
		\framesubtitle{\insertshortsubtitle: Motivation}
		
		\begin{itemize}
			
			\item Classical Seq2Seq models compress the entire source sequence into a single fixed-size vector.
			
			\item \textbf{Problem}:
			\begin{itemize}
				\item Information loss for long sequences.
				\item Difficulty modeling long-range dependencies.
			\end{itemize}
			
			\item \textbf{Intuition}:
			\begin{itemize}
				\item Different target words should focus on different parts of the source sequence.
			\end{itemize}
			
			\item \textbf{Solution}:
			\begin{itemize}
				\item Attention dynamically measures the relevance between elements.
			\end{itemize}
			
		\end{itemize}
		
	\end{frame}
	
	\begin{frame}
		\frametitle{\inserttitle}
		\framesubtitle{\insertshortsubtitle: Plan}
		
		\begin{multicols}{2}
			%	\small
			\tableofcontents
		\end{multicols}
	\end{frame}
	
	\section{Sequence models}
	
	\begin{frame}
		\frametitle{\insertshortsubtitle}
		\framesubtitle{\insertsection}
		
		\begin{itemize}
			\item Sequence models aim to process variable-length structured data.
			\item Goal: learn dependencies across time or token positions.
			\item Challenge: modeling long-range interactions efficiently.
		\end{itemize}
		
	\end{frame}
	
	\subsection{Sequential data}
	
	\begin{subsectionframe}
		\begin{itemize}
			\item \textbf{Definition}:
			\begin{itemize}
				\item Ordered data where each element depends on position in time or sequence.
			\end{itemize}
			
			\item \textbf{Examples}:
			\begin{itemize}
				\item Text: words in a sentence or document.
				\item Speech: audio frames over time.
				\item Time series: financial or sensor data.
			\end{itemize}
			
			\item \textbf{Key property}:
			\begin{itemize}
				\item Order matters; permutation changes meaning.
			\end{itemize}
			
		\end{itemize}
	\end{subsectionframe}
	
	\begin{frame}
		\frametitle{\insertshortsubtitle: \insertsection}
		\framesubtitle{\insertsubsection}
		
		\begin{block}{Sequential data}
			Data where order matters. Observations are not independent.
		\end{block}
		
		\begin{itemize}
			
			\item \textbf{Sequence representation}
			\begin{itemize}
				\item Input shape: $(m, t, d)$ (batch × time steps × feature dimension)
				
				\item Output shape depends on the task:
				\begin{itemize}
					
					\item Sequence labeling: $(m, t, l)$ (batch × time steps × number of classes)
					
					\item Sequence classification: $(m, l)$ (batch × number of classes)
					
					\item Sequence generation: $(m, t', l)$ (batch × generated time steps × vocabulary/classes)
					
				\end{itemize}
			\end{itemize}
			
			\item \textbf{Variable-length sequences}
			\begin{itemize}
				\item Different samples may contain different numbers of time steps.
				\item Padding or masking is often used in mini-batch training.
			\end{itemize}
			
		\end{itemize}
		
	\end{frame}
	
	\subsection{Sequence to sequence model}
	
	\begin{subsectionframe}
		\begin{itemize}
			
			\item \textbf{Architecture}:
			\begin{itemize}
				\item Encoder maps input sequence to a fixed representation.
				\item Decoder generates output sequence step-by-step.
			\end{itemize}
			
			\item \textbf{Training objective}:
			\begin{itemize}
				\item Maximize probability of target sequence given source sequence.
			\end{itemize}
			
		\end{itemize}
	\end{subsectionframe}
	
	\begin{frame}
		\frametitle{\insertshortsubtitle: \insertsection}
		\framesubtitle{\insertsubsection}
		
		\begin{center}
			\hgraphpage[0.8\textwidth]{seq2seq_exp.pdf}
		\end{center}
		
		\begin{itemize}
			\item Encoder--Decoder architecture for sequence transduction.
			\item Suffers from information bottleneck in long sequences.
			\item Decoder relies on fixed-size context vector.
		\end{itemize}
		
		
	\end{frame}
	
	\subsection{Limitations}
	
	\begin{subsectionframe}
		\begin{itemize}
			
			\item \textbf{Long-range dependency problem}:
			\begin{itemize}
				\item Difficulty capturing relationships between distant tokens.
			\end{itemize}
			
			\item \textbf{Information bottleneck}:
			\begin{itemize}
				\item Entire sequence compressed into a fixed-size vector.
			\end{itemize}
			
			\item \textbf{Sequential computation}:
			\begin{itemize}
				\item RNN-based models cannot be fully parallelized.
			\end{itemize}
			
		\end{itemize}
	\end{subsectionframe}
	
	\begin{frame}
		\frametitle{\insertshortsubtitle: \insertsection}
		\framesubtitle{\insertsubsection}
		
		\begin{itemize}
			\item \textbf{Key idea}:
			\begin{itemize}
				\item Instead of compressing all information into one vector, we compute dynamic relevance between tokens.
			\end{itemize}
			
			\item \textbf{Solution}:
			\begin{itemize}
				\item Attention mechanisms allow direct access to all input states.
			\end{itemize}
			
			\item \textbf{Result}:
			\begin{itemize}
				\item Better handling of long-range dependencies and improved interpretability.
			\end{itemize}
		\end{itemize}
		
	\end{frame}
	
	\section{Attention mechanism}
	
	\begin{frame}
		\frametitle{\insertshortsubtitle}
		\framesubtitle{\insertsection}
		
		\begin{center}
			\hgraphpage[0.8\textwidth]{seq2seq_att_exp.pdf}
		\end{center}
		
		\vskip-6pt
		\begin{itemize}
			\item Keep $T$ hidden states from the encoder: $h_i$
			\item For each $s_{t-1}$, calculate its alignment with the hidden states.
			$ e_{t, i} = a(s_{t-1}, h_i) $
			\item Normalize alignment scores into attention weights: $\alpha_{t, i} = \text{softmax}(e_{t, i})$
			\item Calculate the state vector $c_t = \sum_{i=1}^{T} \alpha_{t, i} h_i$
		\end{itemize}
		
	\end{frame}
	
	\subsection{Formalism}
	
	\begin{subsectionframe}
		\begin{itemize}
			\item Learn:
			\begin{itemize}
				\item Keys ($k_i$): representations used to compute attention scores.
				\item Values ($v_i$): representations combined to form the output/context vector.
				\item Queries ($q_j$): representations seeking relevant information.
			\end{itemize}
			\item Calculate the similarity between $q_j$ and $k_i$.
			\item Construct a new representation from this similarity and $v_i$.
			\item Seq2seq attention is a special case of the general attention mechanism.
		\end{itemize}
	\end{subsectionframe}
	
	
	\begin{frame}
		\frametitle{\insertshortsubtitle: \insertsection}
		\framesubtitle{\insertsubsection: General attention}
		
		\hgraphpage[\textwidth]{attention.pdf}
		
	\end{frame}
	
	\begin{frame}
		\frametitle{\insertshortsubtitle: \insertsection}
		\framesubtitle{\insertsubsection: Description (One sample)}
		
		\begin{itemize}
		\item Let $ w_1 \cdots w_{T_s} $ be a source sequence, $x_i = Embedding_{src}(w_i)$
		\[k_i =  x_i W^K \in \mathbb{R}^{d_k} \hskip2cm v_i =  x_i W^V \in \mathbb{R}^{d_v}\]
		
		\item Let $ w'_j $ be the current target word, $y_j = Embedding_{dst}(w'_j)$
		\[q_j =  y_j W^Q \in \mathbb{R}^{d_k}\]

		\item Calculate a similarity between $q_j$ and each key $ k_i $
		\[e_i = a(q_j, k_i)\]
		
		\item Calculate the attention probabilities (softmax can be masked)
		\[\alpha_i = Softmax(e_i) = \frac{exp(e_i)}{\sum_k exp(e_k)}\] 
		
		\item The output is computed as a weighted sum of value vectors:
		\[Attention(q_j, K, V) = \sum_{i=1}^{T_s} \alpha_i v_i\]
		\end{itemize}
		
	\end{frame}
	
	\begin{frame}
		\frametitle{\insertshortsubtitle: \insertsection}
		\framesubtitle{\insertsubsection: Description (Batch)}
		
		\begin{itemize}
			\item $X \in \mathbb{R}^{m \times T_s \times d_x}$ is a source sequence with $m$ samples and $T_s$ time steps
			\[W^K \in \mathbb{R}^{d_x \times d_k} \hskip2cm W^V \in \mathbb{R}^{d_x \times d_v}\]
			\[K =  X W^K \in \mathbb{R}^{m \times T_s \times d_k} \hskip2cm V =  X W^V \in \mathbb{R}^{m \times T_s \times d_v}\]
			
			\item $Y \in \mathbb{R}^{m \times T_t \times d_y}$ is a target sequence with $m$ samples and $T_d$ time steps
			\[W^Q \in \mathbb{R}^{d_y \times d_k} \hskip2cm Q =  Y W^Q \in \mathbb{R}^{m \times T_t \times d_k}\]
			
			\item Calculate similarities between $Q$ and each key $K$
			\[E = a(Q, K) \in \mathbb{R}^{m \times T_t \times T_s} \]
			
			\item Calculate the attention probabilities (softmax can be masked):
			\[A = Softmax(E) \in \mathbb{R}^{m \times T_t \times T_s} \] 
			
			\item The output is computed as a weighted sum of value vectors:
			\[Attention(Q, K, V) = A V \in \mathbb{R}^{m \times T_t \times d_v}\]
		\end{itemize}
		
	\end{frame}
	
	\begin{frame}
		\frametitle{\insertshortsubtitle: \insertsection}
		\framesubtitle{\insertsubsection: Numerical example}
		
		\begin{center}
			\vgraphpage{attention_exp.pdf}
		\end{center}
		
	\end{frame}
	
	\begin{frame}
		\frametitle{\insertshortsubtitle: \insertsection}
		\framesubtitle{\insertsubsection: Gradients (One sample)}
		
		\begin{itemize}
			
			\item Forward relations:
			$
			O = AV, \quad
			A = \text{Softmax}(E), \quad
			E = a(Q,K)$
			
			
			\item Let upstream gradient:
			$
			G = \frac{\partial \mathcal{L}}{\partial O}$
			
			
			\item Gradient w.r.t. values:
			$
			\frac{\partial \mathcal{L}}{\partial V}
			= A^\top G$
			
			
			\item Gradient w.r.t. attention weights:
			$
			\frac{\partial \mathcal{L}}{\partial A}
			= G V^\top$
			
			
			\item Softmax backprop:
			\[
			\frac{\partial \mathcal{L}}{\partial E_{i,j}}
			=
			A_{i,j} \left( H_{i,j} - \sum_{k=1}^{T_s} H_{i,k}A_{i,k} \right)
			\quad \text{where } H = \frac{\partial \mathcal{L}}{\partial A}
			\]
			
			\item Final propagation:
			\[
			\frac{\partial \mathcal{L}}{\partial Q}, \quad
			\frac{\partial \mathcal{L}}{\partial K} \quad
			\text{via chain rule through} \quad
			E = a(Q,K)
			\]
			
		\end{itemize}
		
	\end{frame}
	
	\begin{frame}
		\frametitle{\insertshortsubtitle: \insertsection}
		\framesubtitle{\insertsubsection: Some humor}
		
		\begin{center}
			\vgraphpage{humor/humor-attention.jpeg}
		\end{center}
		
	\end{frame}
	
	\subsection{Attention functions}
	
	\begin{subsectionframe}
		\begin{itemize}
			\item \textbf{Objective}:
			\begin{itemize}
				\item Calculate the similarity between a query and keys
			\end{itemize}
			\item \textbf{Types}:
			\begin{itemize}
				\item \optword{Non parametric}: hard-coded functions.
				\item \optword{Parametric}: learn parameters for the similarity.
			\end{itemize}
			\item For simplicity, we omit the batch dimension.
			\[Q \in \mathbb{R}^{T_t \times d_k}, \, K \in \mathbb{R}^{T_s \times d_k}, \, a(Q, K) \in \mathbb{R}^{T_t \times T_s}, V \in \mathbb{R}^{T_s \times d_v} \]
		\end{itemize}
	\end{subsectionframe}
	
	\begin{frame}
		\frametitle{\insertshortsubtitle: \insertsection}
		\framesubtitle{\insertsubsection: Non parametric (Dot product)}
		
		\begin{itemize}
			\item \textbf{Objective}:
			\begin{itemize}
				\item Measure compatibility between queries and keys.
			\end{itemize}
			
			\item \textbf{Mostly used}:
			\begin{itemize}
				\item Simple attention mechanisms.
				\item Early neural attention models.
				\item When query and key representations share the same vector space.
			\end{itemize}
			
			\item \textbf{Forward pass}:
			\[E = a(Q, K) = Q K^\top\]
			
			\item \textbf{Backward pass (Gradients)}:
			\[
			\frac{\partial \mathcal{L}}{\partial Q}
			=
			\frac{\partial \mathcal{L}}{\partial E} K
			\hskip1cm
			\frac{\partial \mathcal{L}}{\partial K}
			=
			\left(
			\frac{\partial \mathcal{L}}{\partial E}
			\right)^\top Q
			\]
			
		\end{itemize}
		
	\end{frame}
	
	\begin{frame}
		\frametitle{\insertshortsubtitle: \insertsection}
		\framesubtitle{\insertsubsection: Non parametric (Scaled dot product)}
		
		\begin{itemize}
			\item \textbf{Objective}:
			\begin{itemize}
				\item Prevent excessively large dot products when $d_k$ increases.
				\item Stabilize gradients for high-dimensional representations.
			\end{itemize}
			
			\item \textbf{Mostly used}:
			\begin{itemize}
				\item Transformers.
				\item Multi-head attention.
				\item Modern large language models (LLMs).
			\end{itemize}
			
			\item \textbf{Forward pass}:
			\[E = a(Q, K) = \frac{Q K^\top}{\sqrt{d_k}}\]
			
			\item \textbf{Backward pass (Gradients)}:
			\[
			\frac{\partial \mathcal{L}}{\partial Q}
			=
			\frac{1}{\sqrt{d_k}}
			\frac{\partial \mathcal{L}}{\partial E} K
			\hskip1cm
			\frac{\partial \mathcal{L}}{\partial K}
			=
			\frac{1}{\sqrt{d_k}}
			\left(
			\frac{\partial \mathcal{L}}{\partial E}
			\right)^\top Q
			\]
			
		\end{itemize}
		
	\end{frame}
	
	\begin{frame}
		\frametitle{\insertshortsubtitle: \insertsection}
		\framesubtitle{\insertsubsection: Non parametric (Cosine similarity)}
		
		\begin{itemize}
			\item \textbf{Objective}:
			\begin{itemize}
				\item Measure angular similarity between vectors.
				\item Reduce sensitivity to vector magnitude.
			\end{itemize}
			
			\item \textbf{Mostly used}:
			\begin{itemize}
				\item Information retrieval.
				\item Semantic search.
				\item Sentence/document embeddings.
			\end{itemize}
			
			\item \textbf{Forward pass}:
			\[
			e_{ij} = a(q_i, k_j) = \frac{q_i k_j^\top}{||q_i|| \, ||k_j||}
			\]
			
			\item \textbf{Backward pass (Gradients)}:
			\[
			\frac{\partial \mathcal{L}}{\partial q_i}
			=
			\sum_j \delta_{ij} \left( \frac{k_j}{||q_i|| ||k_j||} - \frac{e_{ij}}{||q_i||^2} q_i \right)
			\hskip1cm
			\frac{\partial \mathcal{L}}{\partial k_j}
			=
			\sum_i \delta_{ij} \left( \frac{q_i}{||q_i|| ||k_j||} - \frac{e_{ij}}{||k_j||^2} k_j \right)
			\]
			
		\end{itemize}
		
	\end{frame}
	
	\begin{frame}
		\frametitle{\insertshortsubtitle: \insertsection}
		\framesubtitle{\insertsubsection: Parametric (Additive attention)}
		
		\begin{itemize}
			\item Additive attention (Bahdanau Attention)
			
			\item \textbf{Objective}:
			\begin{itemize}
				\item Learn a nonlinear similarity function.
			\end{itemize}
			
			\item \textbf{Used for}:
			\begin{itemize}
				\item Sequence-to-sequence alignment (e.g., machine translation)
				\item Soft alignment between query and key vectors
				\item Capturing non-linear similarity between representations
			\end{itemize}
			
			\item \textbf{Forward pass}:
			\[e_{ij} = a(q_i, k_j) = v^\top \tanh(W_q q_i + W_k k_j)\]
			
			\item \textbf{Backward pass (Gradients)}:
			\small
			\[
			\frac{\partial \mathcal{L}}{\partial q_i}
			=
			\sum_j \delta_{ij}\, W_q^\top \left( \left(1 - \tanh^2(z_{ij})\right) \odot v \right)
			\hskip.5cm
			\frac{\partial \mathcal{L}}{\partial k_j}
			=
			\sum_i \delta_{ij}\,
			W_k^\top \left( (1 - \tanh^2(z_{ij})) \odot v \right)
			\]
			
			\normalsize where: $z_{ij} = W_q q_i + W_k k_j$ and $ \delta_{ij} =  \frac{\partial \mathcal{L}}{\partial e_{ij}}$
			
		\end{itemize}
	\end{frame}
	
	\subsection{Self-attention}
	
	\begin{subsectionframe}
		\begin{itemize}
			\item \textbf{Objective}:
			\begin{itemize}
				\item Learn to represent a token based on its context (neighbors)
			\end{itemize}
			
			\item \textbf{Used for}:
			\begin{itemize}
				\item Contextual token representation in language models
				\item Sequence modeling (NLP, speech, time series)
				\item Capturing long-range dependencies in sequences
			\end{itemize}
			
			\item \textbf{Characteristics}:
			\begin{itemize}
				\item Queries, keys, and values come from the same sequence.
				\item Enables contextual encoding of tokens.
				\item Forms the basis of transformers.
			\end{itemize}
		\end{itemize}
	\end{subsectionframe}
	
	\begin{frame}
		\frametitle{\insertshortsubtitle: \insertsection}
		\framesubtitle{\insertsubsection: Description}
		
		\begin{itemize}
			%			\item \textbf{GOAL}: Encode an element of a sequence with respect to the entire sequence.
			\item \textbf{Example}: \expword{Encode a word with respect to the sentence to capture its semantics}.
			\item \textbf{Method}:
			\begin{itemize}
				\item Use linear projections to compute queries, keys, and values.
				\item Set the maximum number of elements in a sequence.
				\item Truncate sequences that exceed this number.
				\item Pad those that do not reach this number with a predefined code.
				\item Optionally use multi-head attention to learn multiple representations.
				\item Add positional embeddings to retain positional information in the sequence.
			\end{itemize}
		\end{itemize}
		
	\end{frame}
	
	\begin{frame}
		\frametitle{\insertshortsubtitle: \insertsection}
		\framesubtitle{\insertsubsection: Example}
		
		\begin{center}
			\vgraphpage{self-attention_exp.pdf}
		\end{center}
		
	\end{frame}
	
	\begin{frame}
		\frametitle{\insertshortsubtitle: \insertsection}
		\framesubtitle{\insertsubsection: Some humor}
		
		\hgraphpage{humor/humor-self-attention.jpg}
		
	\end{frame}
	
	\subsection{Multi-Head Attention}
	
	\begin{subsectionframe}
		\begin{itemize}
			\item \textbf{Objective}:
			\begin{itemize}
				\item Learn different representations similar to multiple filters in CNNs.
			\end{itemize}
			
			\item \textbf{Used for}:
			\begin{itemize}
				\item Capturing multiple types of relationships simultaneously (syntax, semantics, position)
				\item Improving representation capacity in Transformer models
				\item Machine translation, language modeling, and sequence understanding tasks
			\end{itemize}
			
			\item \textbf{Advantages}:
			\begin{itemize}
				\item Learn different semantic relationships in parallel.
				\item Attend to information from different representation subspaces.
			\end{itemize}
		\end{itemize}
	\end{subsectionframe}
	
	\begin{frame}
		\frametitle{\insertshortsubtitle: \insertsection}
		\framesubtitle{\insertsubsection: Architecture}
		
		\hgraphpage[\textwidth]{multi_head_att_general.pdf}
		
	\end{frame}
	
	\begin{frame}
		\frametitle{\insertshortsubtitle: \insertsection}
		\framesubtitle{\insertsubsection: Formulas}
		
		\[Q \in \mathbb{R}^{T_t \times d_m}, \, K \in \mathbb{R}^{T_s \times d_m}, \, MultiHead(Q, K, V) \in \mathbb{R}^{T_t \times d_m}, V \in \mathbb{R}^{T_s \times d_m} \]
		
		\[W^Q_i \in \mathbb{R}^{d_m \times d_k}, \,  W^K_i \in \mathbb{R}^{d_m\times d_k}, \, W^V_i \in \mathbb{R}^{d_m \times d_v}, \, W^O \in \mathbb{R}^{h d_v \times d_m}\]
		
		\[head_i = Attention(Q W^Q_i, K W^K_i, V W^V_i)\]
		
		\[MultiHead(Q, K, V) = Concat(head_1, \ldots, head_h) \cdot W^O\]
		
	\end{frame}
	
	\begin{frame}
		\frametitle{\insertshortsubtitle: \insertsection}
		\framesubtitle{\insertsubsection: Some humor}
		
		\begin{center}
			\vgraphpage{humor/humor-multihead.jpg}
		\end{center}
		
	\end{frame}
	
	
	
	\section{Transformers}
	
	\begin{frame}
		\frametitle{\insertshortsubtitle}
		\framesubtitle{\insertsection}
		
		\begin{itemize}
			\item \textbf{Introduced in}:
			\begin{itemize}
				\item \textit{Attention Is All You Need} \cite{2017-vaswani-al}
			\end{itemize}
			
			\item \textbf{Main idea}:
			\begin{itemize}
				\item Replace recurrent computation with attention mechanisms.
			\end{itemize}
			
			\item \textbf{Advantages}:
			\begin{itemize}
				\item Parallel computation.
				\item Better modeling of long-range dependencies.
				\item Scalable training on large corpora.
			\end{itemize}
		\end{itemize}
		
	\end{frame}
	
	\begin{frame}
		\frametitle{\insertshortsubtitle}
		\framesubtitle{\insertsection: Some humor}
		
		\begin{center}
			\vgraphpage{humor/humor-transformer.png}
		\end{center}
		
	\end{frame}
	
	\subsection{Architecture}
	
	\begin{subsectionframe}
		\begin{itemize}
			
			\item \textbf{Core components}:
			\begin{itemize}
				\item Encoder stack and decoder stack.
				\item Multi-head attention and feed-forward layers.
			\end{itemize}
			
			\item \textbf{Key idea}:
			\begin{itemize}
				\item Model relationships between all tokens using self-attention.
			\end{itemize}
			
			\item \textbf{Advantages}:
			\begin{itemize}
				\item Parallelizable training.
				\item Efficient long-range dependency modeling.
			\end{itemize}
			
		\end{itemize}
	\end{subsectionframe}
	
	\begin{frame}
		\frametitle{\insertshortsubtitle: \insertsection}
		\framesubtitle{\insertsubsection: Multi-Head Attention}
		
		\begin{figure}
			\centering
			\hgraphpage[0.9\textwidth]{multi_head_att_.pdf}
			
			%		\vskip-12pt
			\caption{Multi-Head Attention \cite{2017-vaswani-al}}
		\end{figure}
		
	\end{frame}
	
	\begin{frame}
		\frametitle{\insertshortsubtitle: \insertsection}
		\framesubtitle{\insertsubsection: Encoder-decoder}
		
		\begin{minipage}{0.49\textwidth}
			\begin{figure}
				\centering
				\hgraphpage[0.68\textwidth]{transformers_.pdf}
				%			\vskip-8pt
				\caption{Transformers' architecture \cite{2017-vaswani-al}}
			\end{figure}
		\end{minipage}
		\begin{minipage}{0.1\textwidth}
			=
		\end{minipage}
		\begin{minipage}{0.39\textwidth}
			\hgraphpage[\textwidth]{transformers.pdf}
		\end{minipage}
		
	\end{frame}
	
	\subsection{Encoder}
	
	\begin{subsectionframe}
		\begin{itemize}
			
			\item \textbf{Objective}:
			\begin{itemize}
				\item Learn contextual representations of input tokens.
			\end{itemize}
			
			\item \textbf{Components}:
			\begin{itemize}
				\item Multi-head self-attention.
				\item Feed-forward neural networks.
				\item Residual connections and normalization.
			\end{itemize}
			
			\item \textbf{Characteristics}:
			\begin{itemize}
				\item Bidirectional context modeling.
				\item Parallel processing of tokens.
			\end{itemize}
			
		\end{itemize}
	\end{subsectionframe}
	
	\begin{frame}
		\frametitle{\insertshortsubtitle: \insertsection}
		\framesubtitle{\insertsubsection}
		
		\begin{itemize}
			\item Encodes each input element based on its context (the entire input).
			\item Uses multi-head self-attention
			\item Completely bi-directional
			\begin{itemize}
				\item In BiRNN architecture, the encoder represents long-term relation indirectly (using a context).
				\item In Encoders it is directly represented, but the notion of element's position is lost.
			\end{itemize} 
			\item We can introduce position's embedding into the word's initial representation (before being transformed into keys, values and queries).
		\end{itemize}
		
	\end{frame}
	
	\subsection{Decoder}
	
	\begin{subsectionframe}
		\begin{itemize}
			
			\item \textbf{Objective}:
			\begin{itemize}
				\item Generate output tokens sequentially.
			\end{itemize}
			
			\item \textbf{Components}:
			\begin{itemize}
				\item Masked self-attention.
				\item Encoder-decoder attention.
				\item Feed-forward neural networks.
			\end{itemize}
			
			\item \textbf{Characteristics}:
			\begin{itemize}
				\item Autoregressive generation.
				\item Access to encoder representations.
				\item Future tokens are masked during training.
			\end{itemize}
			
		\end{itemize}
	\end{subsectionframe}
	
	\begin{frame}
		\frametitle{\insertshortsubtitle: \insertsection}
		\framesubtitle{\insertsubsection}
		
		\begin{itemize}
			\item \keyword{Autoregressive}: Each token is generated conditioned on previously generated tokens.
			\[p(Y|X) = \prod_{t=1}^{T} p(y_t \mid y_{<t}, X)\]
			
			\item Unlike RNNs:
			\begin{itemize}
				\item Dependencies are modeled directly using masked self-attention.
				\item No recurrent hidden-state propagation is required.
			\end{itemize}
			
			\item Each time an element is generated, it will be added into the input to generate the next element.
			\item To stop generating, a maximum size is used. Also, a special token is used to mark the end of generation.
		\end{itemize}
		
	\end{frame}
	
	\section{Pre-trained models}
	
	\begin{frame}
		\frametitle{\insertshortsubtitle}
		\framesubtitle{\insertsection}
		
		\begin{itemize}
			
			\item Modern transformers are commonly trained in two stages:
			\begin{itemize}
				\item \textbf{Pre-training}: learn general representations from large unlabeled corpora.
				\item \textbf{Fine-tuning}: adapt the model to downstream tasks.
			\end{itemize}
			
			\item This paradigm is known as:
			\begin{itemize}
				\item Self-supervised learning and transfer learning.
			\end{itemize}
			
			\item Different architectures specialize in different objectives:
			\begin{itemize}
				\item Encoder-only.
				\item Decoder-only.
				\item Encoder-decoder (full transformer).
			\end{itemize}
			
		\end{itemize}
		
	\end{frame}
	
	\subsection{Full-transformer}
	
	\begin{subsectionframe}
		\begin{itemize}
			
			\item \textbf{Architecture}:
			\begin{itemize}
				\item Uses both encoder and decoder stacks.
			\end{itemize}
			
			\item \textbf{Objective}:
			\begin{itemize}
				\item Learn sequence-to-sequence transformations.
			\end{itemize}
			
			\item \textbf{Applications}:
			\begin{itemize}
				\item Machine translation.
				\item Summarization.
				\item Text generation and question answering.
			\end{itemize}
			
		\end{itemize}
	\end{subsectionframe}
	
	
	\begin{frame}
		\frametitle{\insertshortsubtitle: \insertsection}
		\framesubtitle{\insertsubsection: T5 \cite{T5}}
		
		\begin{figure}[htbp]
			\hgraphpage[\textwidth]{t5-arch_.pdf}
			\caption{T5 (Text-to-Text Transfer Transformer) training \cite{T5}}
		\end{figure}
		
	\end{frame}
	
	\begin{frame}
		\frametitle{\insertshortsubtitle: \insertsection}
		\framesubtitle{\insertsubsection: BART \cite{bart}}
		
		\begin{figure}[htbp]
			\centering
			\hgraphpage[0.58\textwidth]{bart-arch1_.pdf}
			\hgraphpage[0.4\textwidth]{bart-arch2_.pdf}
			\caption{BART training \cite{bart}}
		\end{figure}
		
	\end{frame}
	
	\subsection{Encoder based}
	
	\begin{subsectionframe}
		\begin{itemize}
			
			\item \textbf{Architecture}:
			\begin{itemize}
				\item Uses only transformer encoders.
			\end{itemize}
			
			\item \textbf{Objective}:
			\begin{itemize}
				\item Learn bidirectional contextual representations.
			\end{itemize}
			
			\item \textbf{Applications}:
			\begin{itemize}
				\item Text classification.
				\item Information retrieval.
				\item Feature extraction and embedding generation.
			\end{itemize}
			
		\end{itemize}
	\end{subsectionframe}
	
	\begin{frame}
		\frametitle{\insertshortsubtitle: \insertsection}
		\framesubtitle{\insertsubsection: BERT \cite{devlin-etal-2019-bert}}
		
		\begin{center}
			\hgraphpage[0.95\textwidth]{bert-arch.pdf}
		\end{center}
		
	\end{frame}
	
	\begin{frame}
		\frametitle{\insertshortsubtitle: \insertsection}
		\framesubtitle{\insertsubsection: ViT \cite{dosovitskiy2021an}}
		
		\vspace{-6pt}
		\begin{figure}[htp!]
			\centering
			\hgraphpage[0.62\textwidth]{vit.pdf}
			%		\vskip-8pt
			\caption{Example of images encoder [modified from \cite{zhang2021dive}]}
		\end{figure}
		
	\end{frame}
	
	\subsection{Decoder based}
	
	\begin{subsectionframe}
		\begin{itemize}
			
			\item \textbf{Architecture}:
			\begin{itemize}
				\item Uses masked transformer decoders only.
			\end{itemize}
			
			\item \textbf{Objective}:
			\begin{itemize}
				\item Predict the next token autoregressively.
			\end{itemize}
			
			\item \textbf{Applications}:
			\begin{itemize}
				\item Text generation.
				\item Dialogue systems.
				\item Large language models (LLMs).
			\end{itemize}
			
		\end{itemize}
	\end{subsectionframe}
	
	\begin{frame}
		\frametitle{\insertshortsubtitle: \insertsection}
		\framesubtitle{\insertsubsection: GPT \cite{radford2018improving}}
		
		\begin{figure}[htbp]
			\hgraphpage[\textwidth]{gpt-arch_.pdf}
			\caption{GPT's architecture and different tasks \cite{radford2018improving}}
		\end{figure}
		
	\end{frame}
	
	\begin{frame}
		\frametitle{\insertshortsubtitle: \insertsection}
		\framesubtitle{\insertsubsection: ImageGPT \cite{chen2020generative}}
		
		\begin{figure}[htbp]
			\hgraphpage[\textwidth]{imagegpt-arch_.pdf}
			\caption{ImageGPT's architecture \cite{chen2020generative}}
		\end{figure}
		
	\end{frame}
	
	\insertbibliography{DL06L-attention}{*}
	
\end{document}